\appendix
\chapter{Secțiuni relevante din cod}\label{ch:anexa-cod}
\pagestyle{fancy}

Anexa reproduce fragmentele de cod la care se face referire în capitolele~\ref{ch:analysis} și \ref{ch:proiectare}, alese pentru că implementează direct rezultatele formulate teoretic. Primele șase secțiuni acoperă cei patru algoritmi din secțiunea~\ref{sec:algoritmi}: generarea unei variante și cele două proceduri pe care le apelează, măsurarea evaziunii împreună cu taxonomia rezultatelor și determinarea importanței determinate prin permutare. Comentariile din cod sunt redate ca atare.

\section{Generarea unei variante perturbate}\label{sec:anexa-generare}

Fragmentul implementează Algoritmul~\ref{alg:generare}. Funcția orchestrează pașii, delegând transformarea propriu-zisă generatorului corespunzător tipului cerut. Propagarea și rezolvarea contorului se produc în interiorul acestuia, iar restaurarea tipurilor încheie procedura. Un nivel în afara intervalului admis produce o eroare, nu o valoare limitată tacit.

\begin{verbatim}
def generate_variant(X, perturbation_type, level, ctx):
    """Genereaza o singura varianta perturbata."""
    if perturbation_type not in PERTURBATION_TYPES:
        raise ValueError(
            f"tip de perturbare necunoscut: {perturbation_type!r}")
    func, n_levels, _ = PERTURBATION_TYPES[perturbation_type]

    # Nivelul 0 e identitatea, folosita drept control: predictiile
    # pe ea trebuie sa reproduca exact baseline-ul.
    if level == 0:
        return identity(X, 0, ctx)
    if not 1 <= level <= n_levels:
        raise ValueError(f"nivel {level} invalid pentru "
                         f"{perturbation_type!r} (1..{n_levels})")

    perturbed, meta = func(X, level, ctx)
    # Propagarea si rezolvarea contorului se fac in `func`; aici se
    # readuc doar tipurile pe care modelul inghetat le asteapta.
    return schema.restore_dtypes(perturbed, X), meta
\end{verbatim}


\section{Propagarea dependențelor}\label{sec:anexa-propagare}

Fragmentul implementează observația centrală din subsecțiunea~\ref{subsec:propagare}. Rapoartele se calculează față de valorile originale, iar caracteristicile derivate se obțin prin înmulțire, nu prin recalculare din formula senzorului. Funcția auxiliară \verb|_safe_ratio| tratează numitorii nuli prin raportul neutru, lăsând caracteristica derivată neschimbată.

\begin{verbatim}
NEUTRAL_RATIO = 1.0

def _safe_ratio(new_values, old_values) -> np.ndarray:
    """Raport element-cu-element, 1.0 unde numitorul e 0."""
    new_arr = np.asarray(new_values, dtype=float)
    old_arr = np.asarray(old_values, dtype=float)
    ratio = np.divide(new_arr, old_arr,
                      out=np.full_like(old_arr, NEUTRAL_RATIO),
                      where=old_arr != 0)
    return np.where(np.isfinite(ratio), ratio, NEUTRAL_RATIO)


def propagate(original, perturbed):
    """Recalculeaza derivatele dupa o schimbare de baza."""
    out = perturbed.copy()

    r_bytes = _safe_ratio(out.sbytes, original.sbytes)
    r_pkts  = _safe_ratio(out.spkts, original.spkts)
    r_dur   = _safe_ratio(out.dur, original.dur)
    # rate ~ (spkts + dpkts - 1) / dur ; dpkts nu se modifica
    r_pktsum = _safe_ratio(out.spkts + out.dpkts - 1,
                           original.spkts + original.dpkts - 1)
    # sinpkt ~ dur / (spkts - 1) : numarul de intervale
    r_gaps = _safe_ratio(out.spkts - 1, original.spkts - 1)

    out["smean"]  = original.smean  * r_bytes / r_pkts
    out["sload"]  = original.sload  * r_bytes / r_dur
    out["rate"]   = original["rate"] * r_pktsum / r_dur
    out["sinpkt"] = original.sinpkt * r_dur / r_gaps
    out["sjit"]   = original.sjit   * r_dur

    # Antrenate cauzal de dilatarea duratei: daca atacatorul
    # intarzie, raspunsurile victimei se intind si ele in timp.
    out["dload"]  = original.dload  / r_dur
    out["dinpkt"] = original.dinpkt * r_dur
    out["djit"]   = original.djit   * r_dur
    return out, stats
\end{verbatim}

\section{Rezolvarea contorului de context}\label{sec:anexa-ct}

Fragmentul materializează cele două politici din subsecțiunea~\ref{subsec:interval}. Politica conservatoare returnează valoarea originală fără a consulta vreo corespondență, iar cea optimistă încearcă, în ordine, tripletul exact observat în date, semnătura traficului legitim pentru aceeași valoare TTL și, în ultimă instanță, o constantă declarată explicit. Statisticile returnate permit raportarea, în fișa rulării, a numărului de rânduri rezolvate pe fiecare cale.

\begin{verbatim}
def resolve_ct_state_ttl(perturbed, policy, lookup, normal_lookup):
    """Determina ct_state_ttl cand sttl/dttl s-au schimbat.

    ct_state_ttl NU e o functie de intervale TTL (sttl=62 -> 2, dar
    sttl=63 -> 0), ci un contor pe fereastra glisanta de 100 de
    conexiuni, imposibil de recalculat din inregistrari izolate.
    """
    if policy not in config.CT_STATE_TTL_POLICIES:
        raise ValueError(
            f"politica ct_state_ttl necunoscuta: {policy!r}")

    if policy == "hold":
        return perturbed["ct_state_ttl"], {...}

    keys = pd.MultiIndex.from_arrays(
        [perturbed["state"], perturbed["sttl"], perturbed["dttl"]])
    exact = lookup.reindex(keys)
    exact_hit = exact.notna().to_numpy()

    from_normal = normal_lookup.reindex(
        perturbed["sttl"]).to_numpy(dtype=float)
    normal_hit = ~np.isnan(from_normal) & ~exact_hit

    values = np.where(exact_hit, exact.to_numpy(dtype=float),
                      np.where(normal_hit, from_normal,
                               config.CT_STATE_TTL_FALLBACK))
    return pd.Series(values, index=perturbed.index), stats
\end{verbatim}

\section{Măsurarea evaziunii pe o variantă}\label{sec:anexa-masurare}

Fragmentul implementează pașii 3--5 din Algoritmul~\ref{alg:masurare}. Varianta nu este citită de pe disc, ci regenerată în memorie, generarea fiind deterministă. Restrângerea la fluxurile eligibile se aplică înaintea oricărei numărători, conform criteriului din subsecțiunea~\ref{subsec:eligibil}.

\begin{verbatim}
def measure_variant(models, references, X, true_labels, ids,
                    variant_type, level, ctx):
    """Masoara o varianta, pentru toate modelele."""
    variant, _ = generators.generate_variant(
        X, variant_type, level, ctx)

    summaries = []
    for model_name, model in models.items():
        reference = references[model_name]
        mask = reference["eligible"]

        predictions = model.predict(variant)

        # Totul se restrange la fluxurile ELIGIBILE: cele pe care
        # modelul le-a semnalat deja pe traficul nemodificat.
        eligible_true = true_labels[mask]
        eligible_pred = np.asarray(predictions)[mask]
        row_outcomes = outcome_utils.classify_outcomes(
            eligible_true, eligible_pred)

        summary = outcome_utils.evasion_summary(row_outcomes)
        summary.update({"model": model_name,
                        "type": variant_type, "level": level})
        summaries.append(summary)
    return summaries
\end{verbatim}


\section{Taxonomia celor trei rezultate}\label{sec:anexa-outcomes}

Fragmentul implementează literal cele trei rezultate din subsecțiunea~\ref{subsec:rezultate}. Ordinea atribuirilor contează: eticheta implicită este confuzia cu alt tip de atac, peste care se suprascriu întâi evaziunea, apoi clasificarea corectă.

\begin{verbatim}
def classify_outcomes(true_labels, predictions) -> np.ndarray:
    """Incadreaza fiecare flux perturbat in una dintre cele
    trei categorii."""
    true_arr = np.asarray(true_labels)
    pred_arr = np.asarray(predictions)

    outcomes = np.full(len(pred_arr),
                  config.OUTCOME_MISCLASSIFIED_ATTACK, dtype=object)
    outcomes[pred_arr == "Normal"] = config.OUTCOME_EVADED
    outcomes[(pred_arr == true_arr)
             & (pred_arr != "Normal")] = config.OUTCOME_CORRECT
    return outcomes
\end{verbatim}

\section{Importanța prin permutare}\label{sec:anexa-importanta-cod}

Fragmentul implementează pașii 1--3 din Algoritmul~\ref{alg:sensibilitate}. Parametrul \verb|donor| materializează corecția discutată în subsecțiunea~\ref{subsec:alg-sensibilitate}: fără el, o analiză restrânsă la o singură clasă ar raporta importanță nulă pentru caracteristici cvasi-constante în acea clasă, deși modelul se poate sprijini decisiv pe ele.

\begin{verbatim}
def permutation_importance(model, X, y_true, columns=None,
                           n_repeats=5, seed=42, donor=None):
    """Importanta prin permutare, pe un model inghetat.

    DONOR: daca analiza e restransa la o clasa in care coloana e
    cvasi-constanta (in Generic, sttl e 254 la 98,8% din randuri),
    amestecarea interna nu schimba nimic si importanta iese fals
    ~0. Cu donor, valorile de inlocuire vin din alta distributie.
    """
    columns = list(columns or X.columns)
    is_attack = y_true != "Normal"
    base_detection, _ = _score(model.predict(X), y_true, is_attack)

    rows = []
    for index, column in enumerate(columns):
        drops = []
        for repeat in range(n_repeats):
            rng = np.random.default_rng(
                seed + 1000 * index + repeat)
            shuffled = X.copy()
            if donor is None:
                shuffled[column] = X[column].to_numpy()[
                    rng.permutation(len(X))]
            else:
                pool = donor[column].to_numpy()
                shuffled[column] = pool[
                    rng.integers(0, len(pool), size=len(X))]
            detection, _ = _score(model.predict(shuffled),
                                  y_true, is_attack)
            drops.append(base_detection - detection)

        rows.append({"feature": column,
                     "detection_drop": float(np.mean(drops)),
                     "detection_drop_std": float(np.std(drops))})

    return pd.DataFrame(rows).sort_values(
        "detection_drop", ascending=False)
\end{verbatim}


\section{Mecanismul de atenție}\label{sec:anexa-atentie}

Fragmentul implementează mecanismul de atenție descris în subsecțiunea~\ref{subsec:atentie}. Atenția este scrisă explicit, fără a folosi componente de nivel înalt din biblioteca de învățare profundă, pentru ca fiecare operație să fie verificabilă. Parametrul \verb|return_attention| expune ponderile, utile pentru interpretabilitate.

\begin{verbatim}
class MultiHeadSelfAttention(nn.Module):
    """Self-attention multi-head, scrisa explicit.

    Pentru fiecare cap h:
        Attention(Q,K,V) = softmax(Q K^T / sqrt(d_head)) V.
    Capetele sunt concatenate si proiectate inapoi in d_token.
    """

    def _split_heads(self, x):
        """[batch, tokens, d] -> [batch, n_heads, tokens, d_head]."""
        batch, tokens, _ = x.shape
        return x.view(batch, tokens,
                      self.n_heads, self.d_head).transpose(1, 2)

    def forward(self, x, return_attention: bool = False):
        q = self._split_heads(self.query(x))
        k = self._split_heads(self.key(x))
        v = self._split_heads(self.value(x))

        scores = q @ k.transpose(-2, -1) / math.sqrt(self.d_head)
        weights = self.dropout(F.softmax(scores, dim=-1))

        context = (weights @ v).transpose(1, 2).reshape(x.shape)
        out = self.projection(context)
        return (out, weights) if return_attention else out
\end{verbatim}

\chapter{Rezultate detaliate}\label{ch:anexa-rezultate}

Anexa conține rezultatele complete din care au fost extrase valorile sintetizate în capitolul~\ref{ch:testare}. Ele sunt reproduse integral pentru ca cifrele raportate în text să poată fi verificate în context.

\section{Rezultatele complete ale măsurării evaziunii}\label{sec:anexa-evaziune}

Tabelul~\ref{tab:anexa-evaziune} prezintă rata de evaziune pentru fiecare dintre cele 28 de variante generate. Capitolul~\ref{ch:testare} raportează doar valorile la intensitate maximă pentru fiecare tip de transformare. Se observă aici comportamentul nemonoton descris în subsecțiunea~\ref{subsec:nemonoton}: pentru transformările TTL, nivelul 2 produce sistematic o rată inferioară nivelului 1, deoarece nivelurile reprezintă valori-țintă distincte, nu intensități crescătoare. Varianta \verb|ttl_both| este marcată ca nerealizabilă și nu trebuie citită alături de celelalte.

\begin{table}[!ht]
    \caption{Rata de evaziune pentru toate cele 28 de variante (\% din fluxurile eligibile)}
    \centering\footnotesize
    \begin{tabular}{|l|c|c|c|c|}
        \hline
        \textbf{Tip} & \textbf{Nivel} & \textbf{RF} & \textbf{XGB} & \textbf{Transformer} \\
        \hline
        combined\_hold & 1 & 0,77 & 10,35 & 5,02 \\
        \hline
        combined\_hold & 2 & 0,80 & 5,17 & 7,12 \\
        \hline
        combined\_hold & 3 & 16,46 & 95,31 & 12,71 \\
        \hline
        combined\_hold & 4 & 31,31 & 96,44 & 17,42 \\
        \hline
        combined\_mimic & 1 & 9,92 & 46,78 & 46,16 \\
        \hline
        combined\_mimic & 2 & 0,81 & 5,15 & 7,35 \\
        \hline
        combined\_mimic & 3 & 66,37 & 96,17 & 66,89 \\
        \hline
        combined\_mimic & 4 & 75,89 & 96,92 & 71,02 \\
        \hline
        connection\_rate & 1 & 0,21 & 0,50 & 0,00 \\
        \hline
        connection\_rate & 2 & 0,20 & 0,56 & 0,00 \\
        \hline
        connection\_rate & 3 & 0,21 & 0,49 & 0,00 \\
        \hline
        connection\_rate & 4 & 0,33 & 0,33 & 0,00 \\
        \hline
        identity & 0 & 0,00 & 0,00 & 0,00 \\
        \hline
        padding & 1 & 0,15 & 1,23 & 0,10 \\
        \hline
        padding & 2 & 0,11 & 0,96 & 0,01 \\
        \hline
        padding & 3 & 0,12 & 0,99 & 0,02 \\
        \hline
        padding & 4 & 0,14 & 1,20 & 0,07 \\
        \hline
        timing & 1 & 0,23 & 0,60 & 0,01 \\
        \hline
        timing & 2 & 0,24 & 0,65 & 0,02 \\
        \hline
        timing & 3 & 0,42 & 0,86 & 0,02 \\
        \hline
        timing & 4 & 0,73 & 1,12 & 0,02 \\
        \hline
        ttl\_both & 1 & 33,02 & 51,07 & 71,07 \\
        \hline
        ttl\_hold & 1 & 0,49 & 7,59 & 3,17 \\
        \hline
        ttl\_hold & 2 & 0,49 & 6,51 & 3,21 \\
        \hline
        ttl\_hold & 3 & 4,44 & 50,21 & 4,04 \\
        \hline
        ttl\_mimic & 1 & 4,46 & 8,74 & 36,98 \\
        \hline
        ttl\_mimic & 2 & 0,50 & 6,51 & 3,47 \\
        \hline
        ttl\_mimic & 3 & 18,42 & 51,07 & 43,15 \\
        \hline
    \end{tabular}
    \label{tab:anexa-evaziune}
\end{table}


\section{Metrici per clasă pentru cele trei modele}\label{sec:anexa-perclasa}

Tabelul~\ref{tab:anexa-perclasa} completează tabelul~\ref{tab:perclasa} din capitolul~\ref{ch:testare}, care raportează doar scorul $F1$. Precizia și sensibilitatea separate arată originea diferenței dintre acuratețea echilibrată și macro-$F1$ discutată în subsecțiunea~\ref{subsec:global}: Transformer obține sensibilități ridicate pe clasele rare (0,83 pe \textit{Backdoor}, 0,88 pe \textit{Shellcode}, 0,80 pe \textit{Worms}), dar cu precizii foarte scăzute, ceea ce ridică prima metrică și o coboară pe a doua.

\begin{table}[!ht]
    \caption{Precizia, sensibilitatea și scorul $F1$ pe clase, pentru cele trei modele}
    \centering\footnotesize
    \begin{tabular}{|l|r|ccc|ccc|ccc|}
        \hline
        & & \multicolumn{3}{c|}{\textbf{RF}} & \multicolumn{3}{c|}{\textbf{XGB}} & \multicolumn{3}{c|}{\textbf{Transformer}} \\
        \textbf{Clasă} & \textbf{Sup.} & P & R & $F1$ & P & R & $F1$ & P & R & $F1$ \\
        \hline
        Normal & 37\,000 & 0,99 & 0,62 & 0,77 & 0,96 & 0,76 & 0,85 & 1,00 & 0,59 & 0,74 \\
        \hline
        Generic & 18\,871 & 1,00 & 0,96 & 0,98 & 1,00 & 0,97 & 0,98 & 1,00 & 0,96 & 0,98 \\
        \hline
        Exploits & 11\,132 & 0,78 & 0,64 & 0,70 & 0,61 & 0,85 & 0,71 & 0,83 & 0,53 & 0,64 \\
        \hline
        Fuzzers & 6\,062 & 0,26 & 0,70 & 0,38 & 0,30 & 0,58 & 0,40 & 0,24 & 0,65 & 0,35 \\
        \hline
        DoS & 4\,089 & 0,43 & 0,15 & 0,23 & 0,57 & 0,12 & 0,20 & 0,33 & 0,20 & 0,25 \\
        \hline
        Reconnaissance & 3\,496 & 0,86 & 0,84 & 0,85 & 0,93 & 0,81 & 0,87 & 0,68 & 0,85 & 0,76 \\
        \hline
        Analysis & 677 & 0,02 & 0,05 & 0,03 & 0,05 & 0,09 & 0,06 & 0,00 & 0,01 & 0,00 \\
        \hline
        Backdoor & 583 & 0,05 & 0,62 & 0,10 & 0,03 & 0,07 & 0,04 & 0,06 & 0,83 & 0,12 \\
        \hline
        Shellcode & 378 & 0,27 & 0,81 & 0,40 & 0,36 & 0,78 & 0,49 & 0,16 & 0,88 & 0,27 \\
        \hline
        Worms & 44 & 0,65 & 0,55 & 0,59 & 0,64 & 0,41 & 0,50 & 0,14 & 0,80 & 0,24 \\
        \hline
    \end{tabular}
    \label{tab:anexa-perclasa}
\end{table}


\section{Importanța completă a caracteristicilor}\label{sec:anexa-importanta}

Tabelul~\ref{tab:anexa-importanta} prezintă clasamentul complet obținut prin permutare, din care capitolul~\ref{ch:testare} reține doar primele trei poziții per model. Sunt listate caracteristicile cu efect nenul pentru cel puțin unul dintre modele. Restul până la cele 42 au produs o scădere a ratei de detecție sub pragul de rezoluție al măsurătorii.

Comparația pe coloane este cea care susține concluzia din subsecțiunea~\ref{subsec:fractiune}: valoarea TTL a sursei domină la XGBoost, în timp ce la Transformer poziția întâi este ocupată de TTL-ul destinației, cu un efect de aproape patru ori mai mare decât al oricărei caracteristici accesibile atacatorului.

\begin{table}[!ht]
    \caption{Importanța prin permutare (scăderea ratei de detecție, exprimată în puncte procentuale). Sunt listate caracteristicile cu efect nenul pentru cel puțin un model.}
    \centering\footnotesize
    \begin{tabular}{|l|r|r|r|}
        \hline
        \textbf{Caracteristică} & \textbf{RF} & \textbf{XGB} & \textbf{Transformer} \\
        \hline
        \texttt{sttl} & 0,993 & 10,995 & 1,286 \\
        \hline
        \texttt{dttl} & 0,533 & 0,044 & 5,064 \\
        \hline
        \texttt{dmean} & 0,087 & 2,107 & 0,063 \\
        \hline
        \texttt{dload} & 0,097 & 1,264 & 0,027 \\
        \hline
        \texttt{sbytes} & -0,092 & 0,952 & -- \\
        \hline
        \texttt{proto} & 0,620 & 0,235 & 0,855 \\
        \hline
        \texttt{smean} & -0,119 & 0,770 & 0,041 \\
        \hline
        \texttt{ct\_dst\_src\_ltm} & 0,170 & 0,722 & 0,189 \\
        \hline
        \texttt{ct\_srv\_src} & -0,017 & 0,635 & 0,022 \\
        \hline
        \texttt{swin} & 0,605 & -- & 0,094 \\
        \hline
        \texttt{sloss} & 0,593 & 0,378 & -0,012 \\
        \hline
        \texttt{djit} & 0,138 & 0,455 & 0,005 \\
        \hline
        \texttt{state} & 0,155 & -0,073 & 0,380 \\
        \hline
        \texttt{spkts} & -0,002 & 0,378 & -- \\
        \hline
        \texttt{ct\_state\_ttl} & 0,310 & 0,136 & 0,298 \\
        \hline
        \texttt{dinpkt} & 0,211 & -0,024 & -- \\
        \hline
        \texttt{ct\_srv\_dst} & -0,070 & 0,194 & 0,170 \\
        \hline
        \texttt{rate} & 0,051 & 0,186 & 0,010 \\
        \hline
        \texttt{dloss} & -0,015 & 0,186 & 0,007 \\
        \hline
        \texttt{sjit} & -0,019 & 0,153 & 0,044 \\
        \hline
        \texttt{dpkts} & 0,073 & 0,114 & 0,007 \\
        \hline
        \texttt{ackdat} & 0,097 & -0,090 & 0,090 \\
        \hline
        \texttt{sload} & -0,065 & 0,094 & -0,002 \\
        \hline
        \texttt{dtcpb} & 0,061 & 0,087 & -0,012 \\
        \hline
        \texttt{dur} & 0,078 & -0,092 & 0,002 \\
        \hline
        \texttt{service} & -0,073 & 0,036 & 0,075 \\
        \hline
        \texttt{dbytes} & 0,073 & -0,659 & -0,002 \\
        \hline
        \texttt{ct\_src\_dport\_ltm} & -0,148 & -0,121 & 0,065 \\
        \hline
        \texttt{ct\_dst\_ltm} & -0,121 & -0,208 & 0,044 \\
        \hline
        \texttt{ct\_flw\_http\_mthd} & 0,017 & 0,024 & -0,002 \\
        \hline
        \texttt{ct\_dst\_sport\_ltm} & 0,019 & -0,068 & -0,005 \\
        \hline
        \texttt{dwin} & 0,012 & -- & 0,019 \\
        \hline
        \texttt{sinpkt} & 0,017 & -0,085 & 0,002 \\
        \hline
        \texttt{trans\_depth} & -0,024 & -0,058 & 0,015 \\
        \hline
        \texttt{tcprtt} & 0,010 & -0,111 & -0,010 \\
        \hline
        \texttt{ct\_src\_ltm} & -0,085 & -0,262 & 0,007 \\
        \hline
        \texttt{response\_body\_len} & -- & -0,012 & 0,005 \\
        \hline
        \texttt{is\_sm\_ips\_ports} & 0,002 & -- & -- \\
        \hline
        \texttt{stcpb} & -0,027 & 0,002 & 0,002 \\
        \hline
        \texttt{synack} & -0,041 & 0,002 & -- \\
        \hline
    \end{tabular}
    \label{tab:anexa-importanta}
\end{table}


\section{Evaziunea per grup de antrenare}\label{sec:anexa-o6}

Tabelele~\ref{tab:anexa-o6-rf}, \ref{tab:anexa-o6-xgb} și \ref{tab:anexa-o6-tr} prezintă ratele complete din care capitolul~\ref{ch:testare} reține doar valoarea din cel mai rău caz. Ele permit două verificări pe care valoarea agregată nu le susține.

Prima privește uniformitatea efectului: la grupul O6, ratele rămân sub $0{,}5\,\%$ pe toate variantele, nu doar pe cea maximă, ceea ce exclude ipoteza că apărarea ar acoperi o singură transformare lăsându-le pe celelalte neatinse.

A doua privește grupul antrenat fără familia TTL. Pe variantele de umplutură, temporizare și reducere a contoarelor, adică cele păstrate în antrenarea lui, valorile sale coincid practic cu ale grupului O6. Pe variantele TTL, excluse, ele revin la nivelul grupului de control pentru două dintre modele. Comparația pe rânduri arată astfel direct unde s-a produs transferul și unde nu.

\begin{table}[!ht]
    \caption{Rata de evaziune a modelului Random Forest pe mulțimea comună, pentru toate variantele realizabile. Valorile sunt procente, mediate peste repetări. Grupul O6-LOO are o singură repetare.}
    \centering\footnotesize
    \begin{tabular}{|l|c|c|c|c|}
        \hline
        \textbf{Variantă} & \textbf{O1} & \textbf{C1} & \textbf{O6} & \textbf{O6-LOO} \\
        \hline
        \texttt{combined\_hold\_L1} & 0,60 & 0,34 & 0,01 & 0,10 \\
        \hline
        \texttt{combined\_hold\_L2} & 0,72 & 0,30 & 0,00 & 0,16 \\
        \hline
        \texttt{combined\_hold\_L3} & 17,23 & 10,14 & 0,00 & 0,72 \\
        \hline
        \texttt{combined\_hold\_L4} & 29,62 & 19,92 & 0,14 & 0,83 \\
        \hline
        \texttt{combined\_mimic\_L1} & 6,78 & 5,82 & 0,00 & 1,17 \\
        \hline
        \texttt{combined\_mimic\_L2} & 0,72 & 0,31 & 0,00 & 0,16 \\
        \hline
        \texttt{combined\_mimic\_L3} & 71,25 & 63,06 & 0,01 & 10,72 \\
        \hline
        \texttt{combined\_mimic\_L4} & 80,32 & 75,51 & 0,18 & 11,84 \\
        \hline
        \texttt{connection\_rate\_L1} & 0,13 & 0,04 & 0,01 & 0,01 \\
        \hline
        \texttt{connection\_rate\_L2} & 0,15 & 0,08 & 0,01 & 0,00 \\
        \hline
        \texttt{connection\_rate\_L3} & 0,19 & 0,11 & 0,00 & 0,00 \\
        \hline
        \texttt{connection\_rate\_L4} & 0,31 & 0,18 & 0,00 & 0,00 \\
        \hline
        \texttt{padding\_L1} & 0,10 & 0,04 & 0,01 & 0,01 \\
        \hline
        \texttt{padding\_L2} & 0,07 & 0,01 & 0,01 & 0,00 \\
        \hline
        \texttt{padding\_L3} & 0,08 & 0,01 & 0,01 & 0,00 \\
        \hline
        \texttt{padding\_L4} & 0,10 & 0,02 & 0,02 & 0,01 \\
        \hline
        \texttt{timing\_L1} & 0,19 & 0,07 & 0,03 & 0,02 \\
        \hline
        \texttt{timing\_L2} & 0,18 & 0,12 & 0,02 & 0,02 \\
        \hline
        \texttt{timing\_L3} & 0,51 & 0,27 & 0,00 & 0,00 \\
        \hline
        \texttt{timing\_L4} & 0,66 & 0,28 & 0,10 & 0,00 \\
        \hline
        \texttt{ttl\_hold\_L1} & 0,37 & 0,18 & 0,00 & 0,03 \\
        \hline
        \texttt{ttl\_hold\_L2} & 0,35 & 0,18 & 0,00 & 0,03 \\
        \hline
        \texttt{ttl\_hold\_L3} & 5,37 & 3,06 & 0,00 & 1,65 \\
        \hline
        \texttt{ttl\_mimic\_L1} & 3,79 & 3,00 & 0,00 & 1,31 \\
        \hline
        \texttt{ttl\_mimic\_L2} & 0,36 & 0,18 & 0,00 & 0,03 \\
        \hline
        \texttt{ttl\_mimic\_L3} & 19,60 & 15,12 & 0,01 & 10,88 \\
        \hline
    \end{tabular}
    \label{tab:anexa-o6-rf}
\end{table}

\begin{table}[!ht]
    \caption{Rata de evaziune a modelului XGBoost pe mulțimea comună, pentru toate variantele realizabile. Valorile sunt procente, mediate peste repetări. Grupul O6-LOO are o singură repetare.}
    \centering\footnotesize
    \begin{tabular}{|l|c|c|c|c|}
        \hline
        \textbf{Variantă} & \textbf{O1} & \textbf{C1} & \textbf{O6} & \textbf{O6-LOO} \\
        \hline
        \texttt{combined\_hold\_L1} & 8,75 & 7,14 & 0,04 & 4,61 \\
        \hline
        \texttt{combined\_hold\_L2} & 4,46 & 3,68 & 0,00 & 2,77 \\
        \hline
        \texttt{combined\_hold\_L3} & 94,72 & 92,79 & 0,00 & 94,02 \\
        \hline
        \texttt{combined\_hold\_L4} & 96,09 & 94,73 & 0,00 & 94,49 \\
        \hline
        \texttt{combined\_mimic\_L1} & 47,04 & 39,56 & 0,02 & 9,30 \\
        \hline
        \texttt{combined\_mimic\_L2} & 4,44 & 3,67 & 0,00 & 2,77 \\
        \hline
        \texttt{combined\_mimic\_L3} & 96,26 & 95,43 & 0,00 & 95,43 \\
        \hline
        \texttt{combined\_mimic\_L4} & 96,90 & 96,24 & 0,00 & 95,68 \\
        \hline
        \texttt{connection\_rate\_L1} & 0,29 & 0,15 & 0,18 & 0,11 \\
        \hline
        \texttt{connection\_rate\_L2} & 0,35 & 0,17 & 0,16 & 0,11 \\
        \hline
        \texttt{connection\_rate\_L3} & 0,34 & 0,11 & 0,05 & 0,01 \\
        \hline
        \texttt{connection\_rate\_L4} & 0,18 & 0,07 & 0,02 & 0,00 \\
        \hline
        \texttt{padding\_L1} & 0,91 & 0,88 & 0,46 & 0,14 \\
        \hline
        \texttt{padding\_L2} & 0,66 & 1,00 & 0,14 & 0,10 \\
        \hline
        \texttt{padding\_L3} & 0,61 & 0,49 & 0,13 & 0,09 \\
        \hline
        \texttt{padding\_L4} & 0,92 & 0,88 & 0,12 & 0,10 \\
        \hline
        \texttt{timing\_L1} & 0,28 & 0,07 & 0,12 & 0,06 \\
        \hline
        \texttt{timing\_L2} & 0,36 & 0,12 & 0,11 & 0,05 \\
        \hline
        \texttt{timing\_L3} & 0,64 & 0,28 & 0,02 & 0,01 \\
        \hline
        \texttt{timing\_L4} & 0,87 & 0,44 & 0,01 & 0,00 \\
        \hline
        \texttt{ttl\_hold\_L1} & 6,47 & 5,73 & 0,00 & 5,95 \\
        \hline
        \texttt{ttl\_hold\_L2} & 5,52 & 4,63 & 0,00 & 5,06 \\
        \hline
        \texttt{ttl\_hold\_L3} & 49,16 & 48,31 & 0,00 & 90,72 \\
        \hline
        \texttt{ttl\_mimic\_L1} & 8,11 & 6,83 & 0,00 & 8,96 \\
        \hline
        \texttt{ttl\_mimic\_L2} & 5,51 & 4,62 & 0,00 & 5,05 \\
        \hline
        \texttt{ttl\_mimic\_L3} & 49,89 & 49,08 & 0,00 & 92,44 \\
        \hline
    \end{tabular}
    \label{tab:anexa-o6-xgb}
\end{table}

\begin{table}[!ht]
    \caption{Rata de evaziune a modelului Transformer pe mulțimea comună, pentru toate variantele realizabile. Valorile sunt procente, mediate peste repetări. Grupul O6-LOO are o singură repetare.}
    \centering\footnotesize
    \begin{tabular}{|l|c|c|c|c|}
        \hline
        \textbf{Variantă} & \textbf{O1} & \textbf{C1} & \textbf{O6} & \textbf{O6-LOO} \\
        \hline
        \texttt{combined\_hold\_L1} & 3,79 & 8,88 & 0,03 & 4,11 \\
        \hline
        \texttt{combined\_hold\_L2} & 4,64 & 12,90 & 0,08 & 6,12 \\
        \hline
        \texttt{combined\_hold\_L3} & 8,26 & 27,48 & 0,08 & 11,40 \\
        \hline
        \texttt{combined\_hold\_L4} & 15,58 & 44,48 & 0,04 & 17,97 \\
        \hline
        \texttt{combined\_mimic\_L1} & 47,86 & 61,96 & 0,09 & 60,91 \\
        \hline
        \texttt{combined\_mimic\_L2} & 4,68 & 12,93 & 0,07 & 6,32 \\
        \hline
        \texttt{combined\_mimic\_L3} & 65,09 & 69,52 & 0,16 & 69,57 \\
        \hline
        \texttt{combined\_mimic\_L4} & 70,92 & 72,52 & 0,09 & 69,65 \\
        \hline
        \texttt{connection\_rate\_L1} & 0,01 & 0,00 & 0,00 & 0,02 \\
        \hline
        \texttt{connection\_rate\_L2} & 0,01 & 0,00 & 0,00 & 0,00 \\
        \hline
        \texttt{connection\_rate\_L3} & 0,00 & 0,00 & 0,00 & 0,00 \\
        \hline
        \texttt{connection\_rate\_L4} & 0,00 & 0,00 & 0,00 & 0,00 \\
        \hline
        \texttt{padding\_L1} & 0,11 & 0,03 & 0,01 & 0,01 \\
        \hline
        \texttt{padding\_L2} & 0,01 & 0,03 & 0,01 & 0,01 \\
        \hline
        \texttt{padding\_L3} & 0,01 & 0,03 & 0,01 & 0,05 \\
        \hline
        \texttt{padding\_L4} & 0,03 & 0,06 & 0,01 & 0,09 \\
        \hline
        \texttt{timing\_L1} & 0,01 & 0,01 & 0,01 & 0,00 \\
        \hline
        \texttt{timing\_L2} & 0,02 & 0,01 & 0,01 & 0,01 \\
        \hline
        \texttt{timing\_L3} & 0,01 & 0,02 & 0,01 & 0,00 \\
        \hline
        \texttt{timing\_L4} & 0,01 & 0,02 & 0,01 & 0,00 \\
        \hline
        \texttt{ttl\_hold\_L1} & 2,98 & 6,45 & 0,02 & 3,75 \\
        \hline
        \texttt{ttl\_hold\_L2} & 3,02 & 6,69 & 0,02 & 3,77 \\
        \hline
        \texttt{ttl\_hold\_L3} & 3,64 & 10,28 & 0,03 & 4,37 \\
        \hline
        \texttt{ttl\_mimic\_L1} & 41,39 & 60,07 & 0,05 & 57,54 \\
        \hline
        \texttt{ttl\_mimic\_L2} & 3,06 & 6,76 & 0,02 & 3,97 \\
        \hline
        \texttt{ttl\_mimic\_L3} & 48,65 & 63,17 & 0,06 & 62,26 \\
        \hline
    \end{tabular}
    \label{tab:anexa-o6-tr}
\end{table}


\section{Configurația experimentală}\label{sec:anexa-config}

Tabelul~\ref{tab:anexa-config} rezumă parametrii înregistrați automat la fiecare rulare completă, conform mecanismului descris în secțiunea~\ref{sec:reproductibilitate}.

\begin{table}[!ht]
    \caption{Configurația cu care au fost produse rezultatele raportate}
    \centering\footnotesize
    \begin{tabular}{|l|l|}
        \hline
        \textbf{Parametru} & \textbf{Valoare} \\
        \hline
        \textit{Seed}-ul generatorului aleator & 42 \\
        \hline
        Partiție & Oficială: 175\,341 / 82\,332 \\
        \hline
        Caracteristici & 42 (39 numerice, 3 categorice) \\
        \hline
        Prag categorii rare & 50 de apariții \\
        \hline
        Random Forest & 200 de arbori, ponderare echilibrată \\
        \hline
        XGBoost & 300 de arbori, adâncime 8, rată 0,1 \\
        \hline
        Transformer & $d = 192$, 3 blocuri, 8 capete, 789\,130 parametri \\
        \hline
        Optimizator & AdamW, rată $10^{-4}$, regularizare $10^{-5}$ \\
        \hline
        Loturi / validare & 512 observații / 15\,\% stratificat \\
        \hline
        Fereastră de selecție & 3 epoci; răbdare 10 epoci \\
        \hline
        Permutări per caracteristică & 5 \\
        \hline
        Predicție & Un singur fir de execuție (determinism) \\
        \hline
        Python & 3.12.10 \\
        \hline
        pandas / numpy & 3.0.5 / 2.4.6 \\
        \hline
        scikit-learn / xgboost & 1.9.0 / 3.4.0 \\
        \hline
        torch & 2.5.1 \\
        \hline
    \end{tabular}
    \label{tab:anexa-config}
\end{table}
